% ECONOMICS_PROBLEM: {"id": "C78-138", "title": "Congestion in matching markets", "jel_code": "C78", "source_id": "OP-0297"}
% !TeX program = xelatex
\documentclass[11pt,a4paper]{article}
\usepackage[margin=23mm]{geometry}
\usepackage{fontspec}
\setmainfont{Latin Modern Roman}
\usepackage{amsmath,amssymb,mathtools}
\usepackage{xcolor,enumitem,booktabs,longtable,array,xurl,hyperref}
\definecolor{muted}{HTML}{53636C}
\hypersetup{colorlinks=true,urlcolor=blue,linkcolor=blue}
\setlength{\parindent}{0pt}
\setlength{\parskip}{5pt}
\newcommand{\R}{\mathbb R}
\newcommand{\E}{\mathbb E}
\newcommand{\N}{\mathbb N}
\newcommand{\ind}{\mathbf 1}
\newcommand{\Var}{\operatorname{Var}}
\newcommand{\Cov}{\operatorname{Cov}}
\newcommand{\supp}{\operatorname{supp}}
\DeclareMathOperator*{\argmin}{arg\,min}
\DeclareMathOperator*{\argmax}{arg\,max}
\newcommand{\entrymeta}[3]{{\sffamily\small\color{muted}\textbf{\texttt{#1}}\quad|\quad #2\par #3\par}}
\newcommand{\Problemref}[1]{\texttt{#1}}
\newcommand{\JELref}[2]{\texttt{#2} (JEL \texttt{#1})}
\begin{document}
\section*{Congestion in matching markets}
\textbf{Problem ID:} \texttt{C78-138}\quad \textbf{JEL:} \texttt{C78}\par
% BEGIN ECONOMICS STATEMENT
\label{problem:OP-0297}
{\sffamily\small\color{muted}Original subject group: Mechanism design and market design\par}
\label{definitions:problem:30007579}
\textbf{Audit classification:} Published research agenda. The verified 2025 paper’s §6 explicitly proposes multiple signaling rounds and richer heterogeneity. The capped-interview BIC frontier is editorial.
\subsection*{Definitions and precise research target}
\textbf{Model and research target.} Let \(A\) be applicants and \(J\) firms in a finite one-to-one matching market. Before interviews, agent \(i\)'s utility for potential partner \(j\) is \(v_j+B_{ij}\); interviewing reveals an additional score \(C_{ij}\). The joint distribution of these scores and public values \(v_j\) is specified. A protocol uses \(r\) signaling rounds, at most \(d\) outgoing signals per agent per round, and at most \(k\) interviews per agent. Signals and interview invitations may depend on earlier messages and revealed scores. Final matches can use only interviewed pairs. An unmatched agent has utility zero. Define interim preferences using revealed total scores for interviewed partners and conditional expected scores for others. A pair blocks if both prefer each other to their final partners under these preferences. Determine the attainable tradeoff between total interviews, signaling rounds, and
\[
 E[\#\{i:i\text{ belongs to an interim blocking pair}\}],
\]
subject to Bayesian incentive compatibility of signaling and an explicitly specified invitation rule. Derive sufficient bounds and matching impossibility examples for richer public-quality heterogeneity and multi-round learning. State how correlations and unbalanced market sizes affect guarantees.

\emph{Definitions and maintained assumptions (editorial unless a source definition is named).} Let \(A,J\) be disjoint finite agent sets; a matching is a set of disjoint applicant--firm edges, and \(M(i)=\varnothing\) means unmatched. Define both sides’ utilities separately, for example \(U_a(j)=v^J_j+B_{aj}+C_{aj}\) and \(U_j(a)=v^A_a+\widetilde B_{ja}+\widetilde C_{ja}\), with unmatched utility zero. Specify each agent’s initial information about these scores and the full correlated prior. An interview reveals the designated pair’s scores according to a declared signal kernel; posterior expected scores of un-interviewed partners incorporate information conveyed by signals and refusals. A pair blocks interim if both strictly prefer it under these posterior utilities, whether or not it interviewed; allowing it to form a final match is a distinct feasibility question. Count each agent involved in any blocking pair once. A multi-round signaling protocol includes invitations, acceptance, score revelation, posterior updating and final matching. Bayesian IC requires truthful signaling to be a best response at each reachable information history, under a specified equilibrium concept and off-path beliefs.
\subsection*{Variants and assumption changes}
\begin{enumerate}
\item \textbf{Round--interview frontier (source direction, editorial criterion).} In balanced random markets with stated i.i.d. scores and public qualities, let market size grow. Determine \((r,d,k)\) requirements for the expected fraction of agents in blocking pairs to tend to zero.
\item \textbf{Heterogeneous and correlated market (source direction, editorial criterion).} Allow multiple quality tiers, unequal side sizes and score correlation. Characterize incentive-compatible signaling and distinguish almost stability from absence of every interim blocking pair.
\end{enumerate}
\subsection*{References and source audit}
\begin{enumerate}
\item §6 Conclusion, printed p.19. Supports the stated research area; exact displayed specification is editorial. \url{https://web.stanford.edu/~alroth/papers/NSF%20Grand%20Challenge.SBE%202020.%20Market%20Design.pdf}
\item §6 Conclusion, printed p.19. Supports the stated research area; exact displayed specification is editorial. \url{https://arxiv.org/pdf/2501.14159}
\end{enumerate}
\begin{enumerate}\raggedright
\item\hypertarget{definitions:source:30007579:1}{} Alvin E. Roth. 2010. \emph{Market Design: Understanding markets well enough to fix them when they’re broken}. §1 Fundamental questions, pp.3–4.
\url{https://web.stanford.edu/~alroth/papers/NSF%20Grand%20Challenge.SBE%202020.%20Market%20Design.pdf}.
\item\hypertarget{definitions:source:30007579:2}{} Maxwell Allman; Itai Ashlagi; Amin Saberi; Sophie H. Yu. 2025. \emph{From signaling to interviews in random matching markets}. §6 Conclusion, p.19, October8,2025 revision.
\url{https://arxiv.org/pdf/2501.14159}.
\end{enumerate}

\subsection*{Common conventions: Source mathematical conventions}

These conventions apply to each entry unless explicitly replaced.

\textbf{Models and identification.} A primitive model \(m\) specifies all probability laws, feasible choices, information, equilibrium and measurement rules used by its target. The admissible class \(\mathcal M\) is fixed independently of outcomes. If \(P_O(m)\) is its observable probability law and \(\tau(m)\) the target, its identified set at observable law \(P\) is
\[
 \mathcal I_\tau(P)=\{\tau(m):m\in\mathcal M,\ P_O(m)=P\}.
\]
Point identification means that this set is a singleton. An empty set signals incompatibility of data and assumptions. Coordinate bounds need not describe the joint set. Bounds are sharp only if every retained target is attainable or arbitrarily approachable by compatible models. Finite bounds require explicit restrictions on unobserved quantities; unrestricted confounding, hidden wealth, extrapolation or arbitrary equilibrium selection can yield uninformative sets.

\textbf{Causal interventions.} A potential outcome \(Y_i(p)\) is the outcome under a complete policy regime \(p\), including timing, financing, information and market spillovers. The target population and units are fixed before treatment. With interference, use the complete assignment vector or a justified exposure mapping; individual no-interference notation alone cannot identify an economy-wide intervention. A difference of mean potential outcomes does not require the cross-world joint distribution. Conditioning on post-treatment migration, survival, enrollment or employment changes the estimand and needs separate justification.

\textbf{Statistical statements.} An estimator, confidence set, forecast or test specifies a sampling law, model class and loss/coverage criterion. Uniform \((1-\alpha)\)-coverage means \(\inf_{m\in\mathcal M}\Pr_m\{\tau(m)\in C_n\}\geq1-\alpha\), or an explicitly asymptotic version. The whole parameter space trivially has coverage, so informativeness requires a stated width, risk or power objective. A forecasting improvement is relative to a named benchmark, information set and evaluation population.

\textbf{Equilibrium and optimization.} An equilibrium comprises feasible plans, optimality, beliefs where needed, budget constraints and all market/resource clearing. The primitives or a stated benchmark define each correspondence \(\mathcal E(p;m)\). Nonemptiness is assumed or proved, never inferred just from compact policy sets. A selected-equilibrium target uses a declared selection rule; a robust target quantifies over all permitted equilibria. Use suprema or infima unless attainment is established. An \(\varepsilon\)-optimal policy is feasible and within \(\varepsilon\) of the finite optimum value. Report an empty feasible set separately.

\textbf{Welfare and accounting.} Normative utility, interpersonal normalization, social weights, numeraire, discounting and the use of public receipts are primitives. Behavioral utility need not coincide with the welfare criterion. Household welfare includes ownership income and rents once; adding profits again to compensating variation generally double-counts them. A monetary equivalent is defined by an explicit transfer and price convention, with monotonicity and range conditions. Average effects, mechanism contributions, transfers and welfare losses are different targets. Infinite sums require convergence and dynamic feasible plans require terminal or no-Ponzi conditions.

\textbf{Source versions.} A working-paper date, revision date and journal publication year are distinguished. A DOI for a journal version is not automatically the DOI of an earlier manuscript. Locators refer to the stated version and printed pages where specified. A metadata-only check does not verify a claim about a full-text conclusion. Every entry's source subsection narrows the provenance claim accordingly.
% END ECONOMICS STATEMENT
\end{document}
